% Sections 7-8 -- discussion and conclusion.  v3, 13 Oct 2026 (FINAL_SUBMISSION_REVIEW_ASTRA: four focused paragraphs; old
% flexible-CG pilot removed; speculative uses reduced to one sentence; design lessons narrowed; conclusion ~180 words).
\section{Discussion}\label{sec:discussion}

\paragraph{What is learned, and what the control shows.} The classical reference receives the same transfers, Galerkin
levels, coarse solution, deflation vectors and output projection as OA-VCB, and the contained fixed-interval Chebyshev member
is within a few per cent of it, so the margins of Section~\ref{sec:accuracy} measure what four learned update stages per
smoother add to one pass through this cycle. That contribution is reproducible across data seeds and training
initialisations, and the ablations locate parts of it under this training protocol: the adapted hierarchy materially improves
the learned cycle (the direct comparison with geometric transfers gives $1.9$ on channels and $4.1$ on barriers), learning
already helps within the residual-linear class ($2.6$--$2.7$ on channels), and the nonlinearity adds a further $3.4$--$3.8$.
The baseline tables of \appref{app:baselines} are supporting context, not evidence for this attribution. The comparison does
not show that the learned pass approximates $A^{-1}$ in any uniform sense.

\paragraph{Why the one-pass advantage does not buy a solver advantage.} Equal update count is a controlled architectural
comparison, not equal work. Additional classical smoothing is cheap on the operator-adapted hierarchy, and with it the margin
falls to $1.8$ on channels and reverses on barriers by the paired metric; three classical passes are more accurate than the
learned pass on both families at about a third of its time (Section~\ref{sec:avt}). Against converged-to-target solvers the gap is one to two
orders of magnitude and grows with the mesh (Section~\ref{sec:cost}); because the network inference alone exceeds a complete
two-grid PCG solve including its set-up, no number of repeated solves amortises training on the tested configurations.
Operator reuse across many right-hand sides and batched GPU execution remain untested and would have to be measured against
equally optimised classical implementations; a correction that is nonlinear in the load would also need a suitable outer
method if used inside an iteration. The operator complexity of the hierarchy ($4.5$ for three levels, $8.6$ for four) would
limit the approach in three dimensions or with more levels unless the prolongations are truncated, as in smoothed aggregation
\cite{vanek1996}, which was not tested.

\paragraph{Loads and floating islands.} The network is positively homogeneous but not linear in the residual. Trained with
one load, it learns corrections for the residuals that this load produces and loses or reverses its advantage on others;
training on a load distribution restores a six-fold margin on that distribution and on smooth loads of both signs, and a
smaller, less uniform one for localised sources (Section~\ref{sec:rhs}). A correction linear in the residual for each fixed
operator would enforce superposition and sign symmetry but would not by itself guarantee accuracy on unseen load directions,
and the tested residual-linear smoothers recover only part of the gain. Indicator deflation improves the tested barrier
models, while substantial error remains on island samples. The tested sequential harmonic enrichment removes only a modest
part of that error; richer coarse spaces, for example contrast-robust spaces built from local spectral problems
\cite{galvis2010,spillane2014geneo,chung2018cem,ma2022} or learned ones \cite{rudikov2025lsino}, are one plausible
continuation, not a remedy established here.

\paragraph{Implications for learned multilevel components.} Four observations bear on how learned components of
multigrid methods are evaluated, beyond the tested implementation. First, a margin measured at one classical smoothing budget
can overstate what learning adds: on the same hierarchy the paired margin fell from $10.2$ to $4.2$ and $1.8$ on channels, and
from $2.1$ to $1.25$ and $0.82$ on barriers, as the Chebyshev degree rose from four to eight and sixteen
(Section~\ref{sec:avt}), so learned smoothers should be compared across that budget. Second, the hierarchy and the learning
interact, and should be varied separately: with geometric transfers the learned cycle was $1.9$ and $4.1$ times less
accurate, and its margin over the classical reference on the same transfers was also smaller ($6.3$--$6.7$ instead of
$9.7$--$10.2$ on channels, $1.48$ instead of $2.13$ on barriers; Table~\ref{tab:transfers}), so in this study the learned
smoothers gained more on the operator-adapted hierarchy than on the geometric one. Third, averages hide where a learned cycle
fails: the error that remained on floating islands after learned smoothing and indicator deflation is visible only when those
samples are reported separately (Section~\ref{sec:islands}). Fourth, a smoother that is nonlinear in the residual is trained
for the residuals of its training loads, so those loads have to sample the loads of the intended use
(Section~\ref{sec:rhs}). Finally, a one-pass accuracy gain at an equal update count does not by itself establish a
computational benefit; that has to be shown at matched accuracy and cost (Section~\ref{sec:cost}).

\paragraph{Scope.} All experiments are two-dimensional with $P_1$ elements; coefficient contrasts are $10^4$ for the
isotropic families, tow coefficients $100$ and $5$ against a unit background for the anisotropic case and a stiffness ratio
of $100$ for the preliminary elasticity case; one family of perforated domains and one polygon are used; models were trained
on meshes of at most 14{,}398 unknowns, and one was applied without retraining to meshes of up to 235{,}774 unknowns, where its
margin shrinks. The learned baselines are adaptations of published architectures to this setting.

\section{Conclusion}\label{sec:conclusion}
We asked what learned smoothing adds to one pass of a multigrid cycle whose transfers, coarse operators, coarsest solve and
deflation vectors are computed from each operator. With four update stages per smoother, the learned pass is about ten times
more accurate in the energy norm than a degree-four Chebyshev-smoothed pass on the same hierarchy for conductive channels and
about twice as accurate for insulating barriers, reproducibly across data seeds and training runs, and the adapted transfers,
deflation, Galerkin output and nonlinearity each contribute. The same experiments bound the result. The margin shrinks quickly
when the classical pass receives more smoothing, and three classical passes are more accurate at a third of the cost; the advantage
depends on the training load; floating islands keep substantial error; and the implementation is much slower than converged
classical solvers at matched accuracy. Learned smoothing is therefore a measurable source of accuracy within a prescribed
multilevel pass, not yet a practical solver. The open challenge is a workload, and an implementation, in which that accuracy
lowers the total cost of reaching a required tolerance.
